\documentclass{article}

\usepackage{arxiv}

\usepackage[utf8]{inputenc} % allow utf-8 input
\usepackage[T1]{fontenc}    % use 8-bit T1 fonts
\usepackage{hyperref}       % hyperlinks
\usepackage{url}            % simple URL typesetting
\usepackage{booktabs}       % professional-quality tables
\usepackage{amsfonts}       % blackboard math symbols
\usepackage{nicefrac}       % compact symbols for 1/2, etc.
\usepackage{microtype}      % microtypography
\usepackage{cleveref}       % smart cross-referencing
\usepackage{lipsum}         % Can be removed after putting your text content
\usepackage{graphicx}
\usepackage{natbib}
\usepackage{doi}

\title{A Motion Capture and Imitation Learning-Based Approach to Robot Control}

% Here you can change the date presented in the paper title
%\date{September 9, 1985}
% Or remove it
%\date{}

\author{ \href{https://orcid.org/0000-0002-8956-179X}{\includegraphics[scale=0.06]{orcid.pdf}\hspace{1mm}Peteris Racinskis}\thanks{Robotics and Machine Perception Laboratory} \thanks{Institute of Electronics and Computer Science, Riga, Latvia} \\
	\texttt{peteris.racinskis@edi.lv} \\
	%% examples of more authors
	\And
	\href{https://orcid.org/0000-0001-5203-3347}{\includegraphics[scale=0.06]{orcid.pdf}\hspace{1mm}Janis Arents}\footnote[1]{test} \footnote[2]{test}\\
	\texttt{janis.arents@edi.lv} \\
	\And
	\href{https://orcid.org/0000-0002-5405-0738}{\includegraphics[scale=0.06]{orcid.pdf}\hspace{1mm}Modris Greitans}\footnote[2]{} \\
	\texttt{modris\_greitans@edi.lv} \\
}

% Uncomment to override  the `A preprint' in the header
%\renewcommand{\headeright}{Technical Report}
%\renewcommand{\undertitle}{Technical Report}
\renewcommand{\shorttitle}{\textit{arXiv} Template}

%%% Add PDF metadata to help others organize their library
%%% Once the PDF is generated, you can check the metadata with
%%% $ pdfinfo template.pdf
\hypersetup{
pdftitle={A Motion Capture and Imitation Learning-Based Approach to Robot Control},
pdfsubject={cs.RO},
pdfauthor={Peteris Racinskis, Janis Arents, Modris Greitans},
pdfkeywords={Imitation Learning, Motion capture, Robotics, Artificial neural networks, RNN},
}

\begin{document}
\maketitle

\begin{abstract}
	Imitation Learning is a discipline of Machine Learning primarily concerned with replicating observed behavior of agents known to perform well on a given task, collected in demonstration data sets. In an industrial robotics context this presents the opportunity to replace explicit programming of behavior with demonstrations of the task to be performed. Motion capture is one of the methods for recording such data. It enables lesser model complexity compared to more indirect observation modalities such as visual data, yet requires additional data pre-processing if signals beyond a time series of effector positions and orientations are relevant to the task at hand. In this paper, an approach for motion capture-based imitation learning and implicit control signal estimation is introduced and evaluated on an object throwing task.
\end{abstract}


% keywords can be removed
\keywords{Imitation Learning \and Motion capture \and Robotics \and Artificial neural networks \and RNN}


\section{Introduction}

Manipulator arms and other types of robots have become ubiquitous in modern industry and their use has been proliferating for decades, yet even now the primary method for programming these devices remains procedural code, hand-crafted by specially trained technicians and engineers. This significantly increases the cost and complexity of commissioning process nodes that utilize industrial robots. Therefore, it is not surprising that attempts to reduce the skill requirements and complexity of robot programming stretch back decades -- with Imitation Learning (IL), or, as it has been otherwise referred to, programming by demonstration, having been proposed and reexamined many times in various technical implementations \citep{muench1994robot,attia2018global}.

As a means of collecting demonstration data, motion capture also has a long history of being studied

Traditionally, theoretical literature in IL has operated on state-action pairs in the context of a Markov Decision Process (MDP)


\paragraph{Motivating use case}
asgasgasgags
asgasgasgags


\section{Related Work}
\label{sec:related}
asgasgasgags
asgasgasgags


asgasgasgagsagaga
asg


\section{Preliminaries}

asfasf

\section{Materials and Methods}
\label{sec:materials}

The overall approach proposed in this paper can be broken down into three main sections -- the collection of observations in the physical environment (\ref{sec:collection}), a pipeline for turning raw observation data into structured demonstrations (\ref{sec:preprocess}) and neural network model implementations (\ref{sec:models}). Two general approaches to system evaluation and design feedback were employed. First, qualitative observations of the generated trajectories were made using spatial visualization in a virtual environment, followed by execution on simulated and real robots (\ref{sec:exec}). Additionally, a series of quantitative metrics have been computed comparing system outputs with training and validation datasets (\ref{sec:eval}).

\subsection{Data collection}
\label{sec:collection}



\subsection{Pre-processing, extraction of implicit control signals}
\label{sec:preprocess}

bla bla bla

\subsection{Models}
\label{sec:models}

bla bla bla

\subsection{Visualization and execution}
\label{sec:exec}


bla bla bla

\subsection{Evaluation metrics}
\label{sec:eval}

bla bla bla


\section{Results}

bla bla bla


\section{Discussion}

bla bla bla



\section{Examples of citations, figures, tables, references}
\label{sec:others}

\subsection{Citations}
Citations use \verb+natbib+. The documentation may be found at
\begin{center}
	\url{http://mirrors.ctan.org/macros/latex/contrib/natbib/natnotes.pdf}
\end{center}

Here is an example usage of the two main commands (\verb+citet+ and \verb+citep+): Some people thought a thing \citep{kour2014real, hadash2018estimate} but other people thought something else \citep{kour2014fast}. Many people have speculated that if we knew exactly why \citet{kour2014fast} thought this\dots


\bibliographystyle{unsrtnat}
\bibliography{references}  %%% Uncomment this line and comment out the ``thebibliography'' section below to use the external .bib file (using bibtex) .



\end{document}
